\subsection{Confession before Christmas}\label{sec:confession}

The voice of Advent is John the Baptist's, and the first thing he preached was
repentance. On the Second Sunday of Advent in Year B the Church hears the
opening of St Mark:

\begin{quote}
John was in the desert, baptizing and preaching the baptism of penance, unto
remission of sins. And there went out to him all the country of Judea and all
they of Jerusalem and were baptized by him in the river of Jordan, confessing
their sins.\par
{\smallskip\noindent\small ---Mark 1:4--5\par}
\end{quote}

\noindent They went out to him. The Advent road leads, among other places, to
the confessional, and the people of Judea went to John before the Lord came.
In Year C, on the Third Sunday, St Luke tells what John said to those who asked
him what their repentance should look like:

\begin{quote}
And the people asked him, saying: What then shall we do? And he answering, said
to them: He that hath two coats, let him give to him that hath none; and he that
hath meat, let him do in like manner. And the publicans also came to be baptized
and said to him: Master, what shall we do? But he said to them: Do nothing more
than that which is appointed you. And the soldiers also asked him, saying: And
what shall we do? And he said to them: Do violence to no man, neither calumniate
any man; and be content with your pay.\par
{\smallskip\noindent\small ---Luke 3:10--14\par}
\end{quote}

The Church's teaching on the sacrament of Penance was set out by the Council of
Trent in its fourteenth session. The acts of the penitent, the Council taught,
are three: contrition, confession and satisfaction. Contrition stands first. It
is sorrow of soul and hatred of the sin committed, together with the purpose of
not sinning again. And the fruit of the sacrament is reconciliation with God,
which, the Council said, is sometimes followed in devout persons who receive
the sacrament with devotion by peace and serenity of conscience, with great
consolation of spirit.

To confess well one must first know one's sins, and the saints have taught
simple ways of looking at a day. St Ignatius of Loyola gave this one in his
Spiritual Exercises, and it serves as well for a season as for a day:

\begin{quote}
It contains in it five Points. First Point. The first Point is to give thanks to
God our Lord for the benefits received. Second Point. The second, to ask grace
to know our sins and cast them out. Third Point. The third, to ask account of
our soul from the hour that we rose up to the present Examen, hour by hour, or
period by period: and first as to thoughts, and then as to words, and then as to
acts \ldots{} Fourth Point. The fourth, to ask pardon of God our Lord for the
faults. Fifth Point. The fifth, to purpose amendment with His grace.\par
{\smallskip\noindent\small ---St Ignatius of Loyola, \emph{Spiritual Exercises}, the General Examen\par}
\end{quote}

\noindent Begin with thanks, as Ignatius does. Then go back over the weeks since
your last confession with John's three questions in mind: what you have kept
that should have been given, where you have taken more than your due, whom you
have treated hardly.

Go early. In the last days before Christmas the lines are long and the priests
are few, and a confession made in the second or third week leaves the last days
free for the joy that John was preparing. An Act of Contrition from the
\emph{Manual of Prayers} of the Third Plenary Council of Baltimore is given in Appendix~\ref{sec:prayers}; learn it by heart if you do not know it,
and teach it to your children before they make their first confession.
